\section{Intermediate FDE Phases Between Sequential Crossings}
\label{sec:intermediate}

Suppose a two-wall transition is not simultaneous. The two unique crossing times then determine an open interval in which exactly one threshold coordinate has already switched to its target side.

The formal development chooses the rational midpoint
\[
  t_m=\frac{t_++t_-}{2}.
\]
If \(t_+<t_-\), Lean proves
\[
  t_+<t_m<t_-.
\]
If \(t_-<t_+\), the symmetric inequalities hold. These midpoint facts require only rational arithmetic.

\subsection{Threshold side before and after a crossing}

The affine support coordinate is strictly monotone whenever its endpoints are distinct. For an increasing straddling path, parameters before the unique crossing remain below threshold and parameters after it lie above threshold. For a decreasing path, the threshold sides reverse. The formal result can be stated without choosing an orientation:

\begin{proposition}[Source side before, target side after]
Let a rational affine support coordinate straddle \(c\) and hit \(c\) uniquely at time \(t_c\). Then
\[
  t<t_c
  \quad\Longrightarrow\quad
  \side_c(\gamma(t))=\side_c(x),
\]
and
\[
  t_c<t
  \quad\Longrightarrow\quad
  \side_c(\gamma(t))=\side_c(y).
\]
\statusverified
\end{proposition}

This is the key bridge from event order to a categorical intermediate state.

\subsection{Generic intermediate vertices}

Let
\[
  s=(s^+,s^-),
  \qquad
  t=(t^+,t^-)
\]
be opposite vertices of the threshold square. Define
\[
  I_+(s,t)=(t^+,s^-)
\]
for positive-first motion and
\[
  I_-(s,t)=(s^+,t^-)
\]
for negative-first motion.

\begin{theorem}[Positive-first midpoint phase]
If both support coordinates straddle and \(t_+<t_-\), then at the rational midpoint
\[
  \text{state}(t_m)=I_+(s,t).
\]
That is, the positive coordinate has already reached the target threshold side while the negative coordinate is still on the source threshold side.
\statusverified
\end{theorem}

\begin{theorem}[Negative-first midpoint phase]
If both support coordinates straddle and \(t_-<t_+\), then
\[
  \text{state}(t_m)=I_-(s,t).
\]
\statusverified
\end{theorem}

For opposite endpoints, both intermediate states are adjacent to source and target:
\[
  \dthr(s,I_\pm(s,t))=1,
  \qquad
  \dthr(I_\pm(s,t),t)=1.
\]
Hence every sequential diagonal displacement decomposes into two one-wall phases.

\subsection{Concrete diagonal table}

The complete verified table is:

\begin{center}
\begin{tabular}{ccc}
\toprule
Diagonal transition & positive wall first & negative wall first\\
\midrule
\(\Nval\to\Bval\) & \(\Tval\) & \(\Fval\)\\
\(\Bval\to\Nval\) & \(\Fval\) & \(\Tval\)\\
\(\Tval\to\Fval\) & \(\Nval\) & \(\Bval\)\\
\(\Fval\to\Tval\) & \(\Bval\) & \(\Nval\)\\
\bottomrule
\end{tabular}
\end{center}

Thus, for example,
\[
  \Nval\to\Bval
\]
can have either phase sequence
\[
  \Nval\to\Tval\to\Bval
\]
or
\[
  \Nval\to\Fval\to\Bval,
\]
and the choice is determined exactly by whether positive or negative support crosses the Lockean wall first.

% Reproducible study figure for the verified affine crossing-order results.
% Threshold c = 0.6. The three affine paths are chosen to illustrate
% positive-first, negative-first, and simultaneous N -> B transitions.
\begin{figure}[t]
\centering
\begin{tikzpicture}[x=10cm,y=10cm,>=Latex]
  % Four threshold cells.
  \fill[black!4] (0,0) rectangle (0.6,0.6);
  \fill[green!8] (0.6,0) rectangle (1,0.6);
  \fill[red!8] (0,0.6) rectangle (0.6,1);
  \fill[blue!5!red!8] (0.6,0.6) rectangle (1,1);

  % Axes and threshold walls.
  \draw[->,thick] (0,0) -- (1.03,0) node[below right] {$P^+$};
  \draw[->,thick] (0,0) -- (0,1.03) node[above left] {$P^-$};
  \draw[dashed,thick] (0.6,0) -- (0.6,1);
  \draw[dashed,thick] (0,0.6) -- (1,0.6);
  \node[below,rotate=90] at (0.615,0.10) {$P^+=c=0.6$};
  \node[above] at (0.13,0.605) {$P^-=c=0.6$};

  % Cell labels.
  \node[align=center,text=black!65] at (0.42,0.18)
    {$\Nval$\\[-1mm]\scriptsize neither / gap};
  \node[align=center,text=green!45!black] at (0.82,0.22)
    {$\Tval$\\[-1mm]\scriptsize true only};
  \node[align=center,text=red!65!black] at (0.22,0.82)
    {$\Fval$\\[-1mm]\scriptsize false only};
  \node[align=center,text=blue!45!red!65!black] at (0.82,0.84)
    {$\Bval$\\[-1mm]\scriptsize both / glut};

  % Positive-first path: (0.2,0.1) -> (0.9,0.7).
  \draw[blue,very thick,->] (0.2,0.1) -- (0.9,0.7);
  \filldraw[blue] (0.2,0.1) circle (0.007);
  \filldraw[blue] (0.9,0.7) circle (0.007);
  \node[blue,below left] at (0.2,0.1) {$x_1$};
  \node[blue,above right] at (0.9,0.7) {$y_1$};
  % t_+ = 4/7, point (0.6, 31/70 ~= 0.443)
  \draw[blue,very thick] (0.6,0.443) node[draw,cross out,minimum size=6pt,inner sep=0pt] {};
  \node[blue,below right,align=left] at (0.61,0.44)
    {\scriptsize $t_+=4/7\approx0.57$\\[-1mm]\scriptsize positive wall first};
  % t_- = 5/6, point (47/60 ~= 0.783, 0.6)
  \draw[blue,very thick] (0.783,0.6) node[draw,cross out,minimum size=6pt,inner sep=0pt] {};
  \node[blue,below right,align=left] at (0.79,0.60)
    {\scriptsize $t_-=5/6\approx0.83$\\[-1mm]\scriptsize negative wall};

  % Negative-first path: (0.1,0.3) -> (0.7,0.9).
  \draw[red!70!yellow,very thick,->] (0.1,0.3) -- (0.7,0.9);
  \filldraw[red!70!yellow] (0.1,0.3) circle (0.007);
  \filldraw[red!70!yellow] (0.7,0.9) circle (0.007);
  \node[red!70!yellow,below left] at (0.1,0.3) {$x_2$};
  \node[red!70!yellow,above] at (0.7,0.9) {$y_2$};
  % t_- = 1/2 at (0.4,0.6)
  \draw[red!70!yellow,very thick] (0.4,0.6) node[draw,cross out,minimum size=6pt,inner sep=0pt] {};
  \node[red!65!black,above left,align=right] at (0.395,0.61)
    {\scriptsize $t_-=1/2$\\[-1mm]\scriptsize negative wall first};
  % t_+ = 5/6 at (0.6,0.8)
  \draw[red!70!yellow,very thick] (0.6,0.8) node[draw,cross out,minimum size=6pt,inner sep=0pt] {};
  \node[red!65!black,above left,align=right] at (0.59,0.81)
    {\scriptsize $t_+=5/6\approx0.83$\\[-1mm]\scriptsize positive wall};

  % Simultaneous path: (0.2,0.2) -> (0.8,0.8), common hit at t=2/3.
  \draw[blue!50!red,densely dotted,very thick,->] (0.2,0.2) -- (0.8,0.8);
  \filldraw[blue!50!red] (0.2,0.2) circle (0.005);
  \filldraw[blue!50!red] (0.8,0.8) circle (0.005);
  \filldraw[blue!50!red] (0.6,0.6) circle (0.010);
  \node[blue!50!red,above left,align=right] at (0.59,0.61)
    {\scriptsize $t_+=t_-=2/3$\\[-1mm]\scriptsize common hit $(c,c)$};

  % Path labels in free regions.
  \node[blue,align=left] at (0.74,0.32)
    {\scriptsize positive-first\\[-1mm]\scriptsize $\Nval\to\Tval\to\Bval$};
  \node[red!65!black,align=left] at (0.28,0.72)
    {\scriptsize negative-first\\[-1mm]\scriptsize $\Nval\to\Fval\to\Bval$};
  \node[blue!50!red,align=left] at (0.38,0.37)
    {\scriptsize simultaneous\\[-1mm]\scriptsize direct diagonal crossing};

  % Unit-square frame.
  \draw[thick] (0,0) rectangle (1,1);
\end{tikzpicture}
\caption{Affine phase geometry for three $\Nval\to\Bval$ support trajectories at threshold $c=0.6$.  Positive-first and negative-first paths necessarily occupy $\Tval$ and $\Fval$, respectively, between their unique wall-crossing times.  The simultaneous path is the exceptional case $t_+=t_-$ and passes through the wall intersection $(c,c)$.  The figure is a visualization of the Lean-verified crossing-order and intermediate-phase theorems; it does not by itself assert that every interpolated support point is already realized by a full admissible model.}
\label{fig:fde-phase-geometry}
\end{figure}


Figure~\ref{fig:fde-phase-geometry} makes the trichotomy explicit for three rational affine trajectories with the same threshold. The two sequential cases pass through different adjacent FDE cells, whereas the simultaneous path hits the unique wall intersection and therefore has no open interval carrying an intermediate adjacent phase.

\begin{theorem}[Two-wall intermediate-phase classification]
Every two-wall conditionalized belief transition admits an affine crossing pair such that either
\begin{enumerate}[label=(\roman*)]
  \item the two crossing times are equal and the support path hits \((c,c)\), or
  \item the positive crossing is earlier and the midpoint has the positive-first intermediate FDE state, or
  \item the negative crossing is earlier and the midpoint has the negative-first intermediate FDE state.
\end{enumerate}
\statusverified
\end{theorem}

\paragraph{Critical semantic boundary.}
The theorem classifies the explicit affine interpolation of the two support masses. It does not yet prove that the midpoint support pair is generated by a complete intermediate 4-PEL model satisfying all model and update constraints. Accordingly, the phrase ``intermediate epistemic phase'' refers here to the thresholded support state along the verified affine support trajectory, not yet to a model-valued dynamical path.
